\documentclass[11pt]{article}
\usepackage[T1]{fontenc}
\usepackage{lmodern}
\usepackage[letterpaper,margin=1in]{geometry}
\usepackage{amsmath,amssymb,amsthm,mathtools,microtype}
\usepackage[hidelinks]{hyperref}
\usepackage[nameinlink,noabbrev]{cleveref}
\usepackage{fancyhdr}
\pagestyle{fancy}\fancyhf{}
\fancyhead[L]{\small Brenier stability under logarithmic second moments}
\fancyhead[R]{\small Research draft v4}
\fancyfoot[C]{\thepage}
\setlength{\headheight}{14pt}
\setlength{\emergencystretch}{2em}
\numberwithin{equation}{section}
\newtheorem{theorem}{Theorem}[section]
\newtheorem{proposition}[theorem]{Proposition}
\newtheorem{lemma}[theorem]{Lemma}
\newtheorem{corollary}[theorem]{Corollary}
\theoremstyle{remark}\newtheorem{remark}[theorem]{Remark}
\crefname{lemma}{lemma}{lemmas}\Crefname{lemma}{Lemma}{Lemmas}
\crefname{proposition}{proposition}{propositions}\Crefname{proposition}{Proposition}{Propositions}
\crefname{corollary}{corollary}{corollaries}\Crefname{corollary}{Corollary}{Corollaries}
\newcommand{\R}{\mathbb R}\newcommand{\E}{\mathbb E}
\newcommand{\Prob}{\mathcal P}\newcommand{\dd}{\,\mathrm d}
\newcommand{\ind}{\mathbf 1}\newcommand{\norm}[1]{\left\lVert#1\right\rVert}
\newcommand{\Var}{\operatorname{Var}}\newcommand{\supp}{\operatorname{supp}}
\newcommand{\id}{\operatorname{id}}
\title{Sharp Gaussian Brenier stability\\under logarithmic second moments}
\author{}\date{Research manuscript 001, version 4\quad 7 October 2026}
\hypersetup{pdftitle={Sharp Gaussian Brenier stability under logarithmic second moments},pdfauthor={},pdfsubject={Sharp logarithmic target-moment exponents and critical finite-moment constants}}
\begin{document}
\maketitle\vspace{-1.5em}
\begin{abstract}
A logarithmic strengthening of the target second moment determines an exact stability exponent for Brenier maps from a fixed Gaussian source. Under the normalized bound $\int (|y|/M)^2[\log(e+(|y|/M)^2)]^\beta\dd\lambda\le1$, $\beta>0$, we prove a homogeneous $L^2$ map estimate with Wasserstein exponent $\min\{1/3,\beta/(\beta+1)\}$. This exponent is optimal in every fixed dimension at least two, including the transition $\beta=1/2$ without a logarithmic loss. At $\beta=0$ there is no uniform modulus in these dimensions. The proof combines the previously established radius-free potential estimate with a minimum-density finite-difference inequality. Integrating signed-square derivative increments directly against the common endpoint weight requires only $L^2$ gradients and produces polynomial source-coordinate weights. A dependence-free logarithmic-moment pairing controls these weights. We also prove uniform estimates for full-support strongly convex $C^{1,1}$ source potentials with a fixed gradient-Lipschitz bound. As a further consequence, the best Gaussian finite-$q$ one-third constant, in its natural homogeneous normalization, has exact order $(q-2)^{-1/6}$ as $q\downarrow2$. The sharpness and endpoint obstructions are multidimensional; in one dimension the map distance equals $W_2$.
\end{abstract}

\paragraph{Status and proof dependency.}
This is a research draft, not a formally verified or externally peer-reviewed paper. No publication-priority claim is made. The all-$\Prob_2$ potential theorem from the public version~3 manuscript \cite{ManuscriptV3} is an explicit input, restated below with its exact hypotheses. The present version gives complete proofs of the new interpolation, logarithmic-moment, sharpness and critical-constant statements; it does not independently re-prove the earlier cell-calculus theorem. Comparison with the full text of \cite{Merigot2026} remains unfinished.

\section{The logarithmic threshold}\label{sec:statements}
Let $\gamma_d$ denote standard Gaussian probability on $\R^d$. For an absolutely continuous source $\rho\in\Prob_2(\R^d)$ and a target $\mu\in\Prob_2(\R^d)$, write $T_\mu$ for the unique quadratic optimal transport map from $\rho$ to $\mu$. It is the gradient of a convex Brenier potential \cite{Brenier1991}. Here $\Prob_q$ denotes probability measures with finite $q$th moment, and
\[
 W_2(\mu,\nu)^2=\inf_{\pi\in\Pi(\mu,\nu)}\int|y-z|^2\dd\pi(y,z).
\]
The Euclidean norm is $|\cdot|$, all logarithms are natural, and $t_+=\max(t,0)$.

For $\beta\ge0$, define $\Phi_\beta(s)=s[\log(e+s)]^\beta$ for $s\ge0$. This function is increasing and is at least $s$. Our target condition, with common scale $M>0$, is
\begin{equation}\label{eq:log-class}
 \int\Phi_\beta\left(\frac{|y|^2}{M^2}\right)\dd\lambda(y)\le1
 \qquad(\lambda=\mu,\nu).
\end{equation}
It implies both target second moments are at most $M^2$.

\begin{theorem}[Sharp Gaussian logarithmic-moment estimate]\label{thm:log}
Let $d\ge1$, $\beta>0$, and put
\begin{align}
 \eta&=\min\{1,2\beta\},&
 \alpha&=\frac\eta{\eta+2}=\min\left\{\frac13,\frac\beta{\beta+1}\right\},\label{eq:eta-alpha}\\
 K_\eta&=6\cdot2^\eta+2,&
 B_\beta&=2+8^\beta+\sqrt2(8\beta/e)^\beta,\label{eq:constants}\\
 C_{d,\beta}&=\bigl(24d+2dK_\eta B_\beta\bigr)^{1/2}.&&\notag
\end{align}
If $\mu,\nu$ satisfy \eqref{eq:log-class}, their Brenier maps from $\gamma_d$ obey
\begin{equation}\label{eq:main-log}
 \norm{T_\mu-T_\nu}_{L^2(\gamma_d)}
 \le C_{d,\beta}M^{1-\alpha}W_2(\mu,\nu)^\alpha.
\end{equation}
Moments may instead be measured about any common center $z\in\R^d$.

For each fixed $d\ge2$, the exponent $\alpha$ cannot be increased on the normalized class \eqref{eq:log-class} with $M=1$. At $\beta=0$, that class has no uniform map-continuity modulus. In dimension one,
\begin{equation}\label{eq:one-dimensional}
 \norm{T_\mu-T_\nu}_{L^2(\gamma_1)}=W_2(\mu,\nu),
\end{equation}
including the second-moment endpoint.
\end{theorem}

The transition at $\beta=1/2$ is included: its exponent is exactly $1/3$. The constants displayed here are admissible, not asserted optimal. For $\beta>0$ the condition already implies uniform integrability of the second moments, since
\[
 \int_{|y|>R}|y|^2\dd\lambda(y)
 \le\frac{M^2}{[\log(e+(R/M)^2)]^\beta}.
\]
Thus qualitative uniform continuity also follows from classical transport stability and $W_2$ compactness. The issue addressed here is the explicit sharp power law and its endpoint.

\subsection{A smooth-source extension}
The same exponent holds for full-support uniformly strongly log-concave sources with a controlled gradient of their logarithmic density.

\begin{theorem}[Uniform smooth-source estimate]\label{thm:smooth}
Let $V:\R^d\to\R$ be continuously differentiable, $\kappa$-strongly convex, and have globally $\Lambda$-Lipschitz gradient, where $0<\kappa\le\Lambda<\infty$. Let $\rho=Z^{-1}e^{-V}\dd x$. For $\beta>0$ set $\eta,\alpha,K_\eta$ as in \eqref{eq:eta-alpha}--\eqref{eq:constants} and define
\begin{align}
 B_{d,\kappa,\Lambda,\beta}
 &=2+(8/\kappa)^\beta
     +(8\beta/(\kappa e))^\beta(2\Lambda/\kappa)^{d/2},\label{eq:smooth-B}\\
 C_{d,\kappa,\Lambda,\beta}
 &=\max\left\{2,
 \left(\frac{24d}{\kappa}
       +2dK_\eta\Lambda^\eta B_{d,\kappa,\Lambda,\beta}\right)^{1/2}\right\}.
 \label{eq:smooth-C}
\end{align}
For targets satisfying \eqref{eq:log-class},
\begin{equation}\label{eq:smooth-main}
 \norm{T_\mu-T_\nu}_{L^2(\rho)}
 \le C_{d,\kappa,\Lambda,\beta}M^{1-\alpha}W_2(\mu,\nu)^\alpha.
\end{equation}
The constant is uniform in $V$ for fixed $d,\kappa,\Lambda,\beta$, and common-center moments may be used. If $d\ge2$, the exponent is sharp over this source class; individual sharpness for every non-Gaussian member is not asserted. Hard convex boundaries are excluded from this theorem.
\end{theorem}

\subsection{The critical finite-moment constant}
The same interpolation improves the dependence on the target-moment order in the Gaussian one-third estimate.

\begin{theorem}[Sharp finite-$q$ constant order]\label{thm:q}
For $q>2$, put
\[
 p=\frac q{q-2},\qquad D_q=\norm{1+|G|}_{L^p(\gamma_1)},\qquad G\sim\gamma_1,
\]
and, for $\mu,\nu\in\Prob_q(\R^d)$, let
\[
 M_q=\left(\int|y|^q\dd\mu(y)\right)^{1/q}
     +\left(\int|y|^q\dd\nu(y)\right)^{1/q}.
\]
Then for the standard Gaussian source in every $d\ge1$,
\begin{equation}\label{eq:q-upper}
 \norm{T_\mu-T_\nu}_{L^2(\gamma_d)}
 \le\sqrt{38d}\,D_q^{1/3}M_q^{2/3}W_2(\mu,\nu)^{1/3}.
\end{equation}
Let $C_*(d,q)$ be the smallest constant in the same inequality when the factor $\sqrt{38d}D_q^{1/3}$ is replaced by a constant uniform over all such targets. For $d\ge2$, there are explicit $c_0>0$ and $A_d<\infty$, independent of $q$, with
\begin{equation}\label{eq:q-sharp-order}
 c_0p^{1/6}\le C_*(d,q)\le A_dp^{1/6}.
\end{equation}
In particular $C_*(d,q)\asymp_d(q-2)^{-1/6}$ as $q\downarrow2$. For $d=1$, $C_*(1,q)=1$.
\end{theorem}

The last theorem concerns the best constant in the specified homogeneous normalization, rather than a different transport exponent. Explicit choices for $c_0,A_d$ appear in \Cref{sec:q}.

\subsection{Provenance and related estimates}
The potential estimate used here is \cite[Theorem~1.1]{ManuscriptV3}. That manuscript develops a strongly log-concave extension of the finite-cell and moving-site framework of the OpenAI uniform-source preprint \cite{OpenAI2026}. The earlier Gaussian three-atom construction is used below with full details. The new interpolation keeps signed derivative increments under the minimum of the two endpoint densities; it does not estimate an exponential density ratio against the original weight.

Potential-to-gradient interpolation and unbounded-target moment estimates have important precedents. Delalande and M\'erigot \cite[Proposition~4.1 and Corollary~4.4]{DelalandeMerigot2023} treat bounded-gradient interpolation and finite target moments for compact convex sources with two-sided density bounds. For moment order $p\ge4$, $p>d$, their map exponent is $p/(6p+16d)$ in $W_1$. Letrouit and M\'erigot \cite[Theorem~1.4 and Remark~1.5]{LetrouitMerigot2026} treat strongly log-concave sources and bounded logarithmic perturbations with compact targets; their unbounded-target case is left untreated there. Mischler and Trevisan \cite{MischlerTrevisan2026} establish related potential and map estimates for general power costs and discuss weighted interpolation. These results have different source or target hypotheses. Neither a metric exponent nor the absence of a statement in the compared results is evidence of priority. A full comparison with \cite{Merigot2026} remains unfinished.

\section{The inherited potential theorem}\label{sec:input}
We state exactly the result from version~3 that is required below. Its full public proof, including nonsmooth source conditioning, finite-target approximation and completion to all $\Prob_2$ targets, is in the cited manuscript.

\begin{theorem}[Potential stability input, {\cite[Theorem~1.1]{ManuscriptV3}}]\label{thm:input}
Let $\rho=r\dd x$ be full-dimensional with
\[
 r(x)=Z^{-1}\exp(-\kappa|x|^2/2-W(x)),\qquad \kappa>0,
\]
where $W:\R^d\to(-\infty,+\infty]$ is proper, lower-semicontinuous and convex. For every $\mu\in\Prob_2(\R^d)$, its Brenier potential has a unique centered representative $U_\mu\in L^2(\rho)$, finite and convex on the interior of the source support. For every $\mu,\nu\in\Prob_2(\R^d)$,
\begin{equation}\label{eq:potential-input}
 \norm{U_\mu-U_\nu}_{L^2(\rho)}^2
 \le\frac2\kappa W_2(\mu,\nu)^2.
\end{equation}
\end{theorem}

This is an established input within the cited proof chain, not a new conditional hypothesis on the targets. For the Gaussian and smooth sources of the present paper the support is all of $\R^d$, so the representatives are finite convex functions everywhere. The Gaussian has $W=0$, $\kappa=1$. Under \Cref{thm:smooth}, $W=V-\kappa|x|^2/2$ is finite convex. The source has a finite second moment in both cases. No bounded-gradient conclusion is inferred for unbounded targets.

We also use the standard optimality of a convex gradient to its own law. For finite maxima of affine functions, it follows directly by integrating the supporting affine inequalities: the gradient plan maximizes $\int x\cdot y$, hence minimizes quadratic cost, and equality forces its unique maximizing slope almost everywhere. This verifies every explicit map in the sharpness examples below.

\section{A minimum-density finite-difference inequality}\label{sec:minimum}
The key point is to preserve the common endpoint weight when integrating the signed-square derivative increment. The following form does not assume log-concavity of the density.

\begin{lemma}[General minimum-weight inequality]\label[lemma]{lem:minimum}
Let $\rho=r\dd x$ be a probability on $\R^d$ with $0<r<\infty$ almost everywhere. Let $u,v$ be finite convex functions on $\R^d$, with gradients in $L^2(\rho)$ and $f=u-v-a\in L^2(\rho)$ for some $a\in\R$. For $h>0$ and each coordinate $j$, define
\begin{align*}
 m_{j,h}(x)&=\min\{r(x),r(x+he_j)\},\\
 A_{j,h}(x)&=\frac{|m_{j,h}(x-he_j)-m_{j,h}(x)|}{r(x)},\qquad
 B_{j,h}(x)=1-\frac{m_{j,h}(x)}{r(x)}.
\end{align*}
Then $0\le A_{j,h},B_{j,h}\le1$ almost everywhere and
\begin{equation}\label{eq:general-minimum}
 \begin{split}
 \norm{\nabla u-\nabla v}_{L^2(\rho)}^2
 &\le\frac{12d}{h^2}\norm f_{L^2(\rho)}^2\\
 &\quad+\sum_{j=1}^d\int(6A_{j,h}+2B_{j,h})
                (|\partial_j u|^2+|\partial_j v|^2)\dd\rho.
 \end{split}
\end{equation}
\end{lemma}
\begin{proof}
Fix $j,h$ and abbreviate $m=m_{j,h}$. For this fixed displacement, discard the null differentiability sets and their translates. Convexity gives
\[
 \partial_j u(x)\le D_hu(x):=\frac{u(x+he_j)-u(x)}h
          \le\partial_j u(x+he_j).
\]
Set $a_u=\partial_j u$, $g_u=a_u|a_u|$ and $\delta_u=D_hu-a_u\ge0$. The elementary inequality
\begin{equation}\label{eq:signed-square}
 (b-a)^2\le2(b|b|-a|a|)\qquad(a\le b)
\end{equation}
follows by factoring the difference of squares when the signs agree, and by $(|a|+|b|)^2\le2(a^2+b^2)$ when they differ. Consequently
\[
 \delta_u(x)^2\le2\bigl(g_u(x+he_j)-g_u(x)\bigr),
\]
with a nonnegative bracket.

Both terms are absolutely integrable under the minimum weight using only the assumed $L^2$ gradient:
\begin{align*}
 \int|g_u(x)|m(x)\dd x&\le\int|a_u|^2r\dd x<\infty,\\
 \int|g_u(x+he_j)|m(x)\dd x
 &=\int|g_u(y)|m(y-he_j)\dd y\le\int|a_u|^2r\dd y<\infty.
\end{align*}
The second bound uses $m(y-he_j)=\min(r(y-he_j),r(y))\le r(y)$. Signed subtraction and a change of variable are therefore legitimate. They give
\begin{align}
 \int\delta_u^2m\dd x
 &\le2\int g_u(y)\bigl(m(y-he_j)-m(y)\bigr)\dd y\notag\\
 &\le2\int|a_u|^2 A_{j,h}\dd\rho.\label{eq:minimum-error}
\end{align}
The first signed integral is nonnegative by its preceding representation; its integrand after translation need not be pointwise nonnegative. Taking an absolute-value upper bound is valid. The same calculation holds for $v$.

The common weight controls both finite-difference endpoints:
\begin{align}
 \int|D_hf|^2m\dd x
 &\le\frac2{h^2}\int\bigl(|f(x+he_j)|^2+|f(x)|^2\bigr)m(x)\dd x\notag\\
 &\le\frac4{h^2}\norm f_{L^2(\rho)}^2.\label{eq:fd-bound}
\end{align}
No translated $L^2(\rho)$ assumption on $f$ is used. Since
$a_u-a_v=D_hf-\delta_u+\delta_v$, the three-term square inequality and \eqref{eq:minimum-error}--\eqref{eq:fd-bound} yield
\[
 \int|a_u-a_v|^2m\dd x
 \le\frac{12}{h^2}\norm f_2^2
          +6\int A_{j,h}(|a_u|^2+|a_v|^2)\dd\rho.
\]
On the remaining weight, $|a_u-a_v|^2\le2(|a_u|^2+|a_v|^2)$ gives the term with $2B_{j,h}$. Summing in $j$ proves \eqref{eq:general-minimum}. The bounds $A_{j,h},B_{j,h}\le1$ follow because both minimum densities in $A_{j,h}$ lie between zero and $r(x)$.
\end{proof}

\begin{corollary}[Gaussian polynomial-weight interpolation]\label[corollary]{cor:gaussian-interpolation}
Under the assumptions of \Cref{lem:minimum} with $\rho=\gamma_d$, for $0<\eta\le1$ and $0<h\le1$,
\begin{equation}\label{eq:gaussian-interpolation}
 \begin{split}
 \norm{\nabla u-\nabla v}_2^2
 &\le12dh^{-2}\norm{u-v-a}_2^2\\
 &\quad+K_\eta h^\eta\sum_j\int
 (|\partial_j u|^2+|\partial_j v|^2)(1+|x_j|)^\eta\dd\gamma_d.
 \end{split}
\end{equation}
The last integrals may initially be infinite; the inequality remains valid in the extended-real sense.
\end{corollary}
\begin{proof}
Put $F(s)=\min(1,e^s)$. This function takes values in $(0,1]$ and is globally $1$-Lipschitz. Gaussian translation gives
\[
 \frac{m_{j,h}(x-he_j)-m_{j,h}(x)}{r(x)}
 =F(hx_j-h^2/2)-F(-hx_j-h^2/2).
\]
Hence $A_{j,h}\le\min(2h|x_j|,1)$. Also
\[
 B_{j,h}=(1-e^{-hx_j-h^2/2})_+
 \le\min\{h(|x_j|+\tfrac12),1\}.
\]
For $t\ge0$ and $0<\eta\le1$, $\min(t,1)\le t^\eta$. Therefore
\[
 6A_{j,h}+2B_{j,h}
 \le h^\eta\bigl(6\cdot2^\eta|x_j|^\eta
                     +2(|x_j|+\tfrac12)^\eta\bigr)
 \le K_\eta h^\eta(1+|x_j|)^\eta.
\]
Insert this into \eqref{eq:general-minimum}.
\end{proof}

\section{A dependence-free logarithmic pairing}\label{sec:pairing}
The map coordinate and the source coordinate in the next application are generally dependent. No independence is needed.

\begin{lemma}[Logarithmic moment pairing]\label[lemma]{lem:pairing}
On a probability space, let $A,S\ge0$ be arbitrary jointly distributed variables. Assume $\beta>0$, $a>0$ and
\[
 \E\Phi_\beta(A)\le1,\qquad \E e^{aS^2}\le H<\infty.
\]
If $0<\eta\le\min(1,2\beta)$, then
\begin{equation}\label{eq:pairing}
 \E[A(1+S)^\eta]
 \le2+(2/a)^\beta+(2\beta/(ae))^\beta H.
\end{equation}
\end{lemma}
\begin{proof}
Since $\Phi_\beta(A)\ge A$, $\E A\le1$. For $s\ge0$,
$(1+s)^\eta\le1+s^\eta\le2+s^{2\beta}$. On the event
$S^2\le(2/a)\log(e+A)$,
\[
 AS^{2\beta}\le(2/a)^\beta\Phi_\beta(A).
\]
On the complementary event, $A<e^{aS^2/2}$, and
\[
 AS^{2\beta}\le e^{aS^2/2}S^{2\beta}
 \le(2\beta/(ae))^\beta e^{aS^2}.
\]
The last inequality follows by maximizing $x^\beta e^{-ax/2}$ over $x\ge0$, whose maximum is $(2\beta/(ae))^\beta$. Taking expectations proves \eqref{eq:pairing}.
\end{proof}

For $S=|G|$, $G\sim\gamma_1$, take $a=1/4$ and $H=\sqrt2$, since $\E e^{G^2/4}=\sqrt2$. Then \eqref{eq:pairing} gives exactly $B_\beta$ in \eqref{eq:constants}.

\begin{proof}[Upper estimate in \Cref{thm:log}]
The logarithmic condition puts both targets in $\Prob_2$. Let $u=U_\mu$, $v=U_\nu$ be the centered potentials from \Cref{thm:input}, finite on all of $\R^d$. Set $w=W_2(\mu,\nu)$. Their difference satisfies $\norm{u-v}_2^2\le2w^2$.
For each $j$, apply \Cref{lem:pairing} with
\[
 A=|\partial_j u|^2/M^2,\qquad S=|x_j|.
\]
The source coordinate is standard Gaussian. Monotonicity of $\Phi_\beta$, the inequality $|\partial_j u|\le|\nabla u|$, and $(\nabla u)_\#\gamma_d=\mu$ verify its logarithmic hypothesis. The same holds for $v$. Hence
\[
 \sum_j\int(|\partial_j u|^2+|\partial_j v|^2)
               (1+|x_j|)^\eta\dd\gamma_d\le2dB_\beta M^2.
\]
By \Cref{cor:gaussian-interpolation},
\begin{equation}\label{eq:log-balance}
 \norm{T_\mu-T_\nu}_2^2
 \le24d w^2h^{-2}+2dK_\eta B_\beta M^2h^\eta.
\end{equation}
If $0<w\le M$, choose $h=(w/M)^{2/(\eta+2)}\in(0,1]$. Both terms have scale $M^{4/(\eta+2)}w^{2\eta/(\eta+2)}$, and taking square roots proves \eqref{eq:main-log}. If $w>M$, the elementary bound $\norm{T_\mu-T_\nu}_2\le2M$ proves the same conclusion since $C_{d,\beta}\ge2$. If $w=0$, uniqueness gives equality of the maps. The scale $M$ is strictly positive, so no division by zero is involved.

For a common center $z$, subtract $z\cdot x$ from both convex potentials. Their difference and the map difference remain unchanged, and the gradient laws are the common translates of the targets. The same argument therefore proves the centered-moment version.
\end{proof}

\section{Sharpness and the second-moment endpoint}\label{sec:sharpness}
We first establish a reusable two-atom calculation. Its radius will then be chosen to fit the logarithmic normalization or a fixed power moment.

\begin{lemma}[Rare rotated halfspaces]\label[lemma]{lem:halfspaces}
Let $d\ge2$, $t\ge1$, and let $G,Z$ be independent standard normals. Put
\[
 m_t=\Pr(G>t),\qquad n_t=\cos(1/t)e_1+\sin(1/t)e_2,
 \qquad\delta_t=|e_1-n_t|.
\]
For any $R>0$, define
\[
 \mu=(1-m_t)\delta_0+m_t\delta_{Re_1},\qquad
 \nu=(1-m_t)\delta_0+m_t\delta_{Rn_t}.
\]
Their Gaussian Brenier maps are $Re_1\ind_{A_t}$ and $Rn_t\ind_{B_t}$, where
$A_t=\{x_1>t\}$ and $B_t=\{x\cdot n_t>t\}$. If $D_t=\gamma_d(A_t\mathbin\triangle B_t)$, then
\begin{align}
 W_2(\mu,\nu)^2&=m_tR^2\delta_t^2,\label{eq:halfspace-W}\\
 \norm{T_\mu-T_\nu}_2^2
 &=R^2\bigl[D_t+\delta_t^2(m_t-D_t/2)\bigr],\label{eq:halfspace-map}\\
 c\,m_t&\le D_t\le2m_t,
 \qquad c=e^{-3/2}\Pr(Z<-2)>0.\label{eq:halfspace-D}
\end{align}
Also
\begin{equation}\label{eq:mills}
 e^{-3/2}\frac{\varphi(t)}t\le m_t\le\frac{\varphi(t)}t,
 \qquad \delta_t=2\sin(1/(2t))\sim t^{-1},
\end{equation}
where $\varphi$ is the standard normal density.
\end{lemma}
\begin{proof}
The maps are gradients of $R(x_1-t)_+$ and $R(x\cdot n_t-t)_+$. Rotational invariance gives both rare cells mass $m_t$; convexity and the supporting-inequality argument identify their optimal maps.
Since $m_t<1/2$, a two-by-two target coupling is specified by the mass $s\in[0,m_t]$ sent between the two nonzero atoms. Its squared cost is
\[
 R^2\bigl(2(m_t-s)+s\delta_t^2\bigr)
 =2R^2\bigl(m_t-s\cos(1/t)\bigr).
\]
This is minimized at $s=m_t$ because $\cos(1/t)>0$, proving \eqref{eq:halfspace-W}. On the symmetric difference the squared map discrepancy is $R^2$; on the overlap it is $R^2\delta_t^2$. The overlap mass is $m_t-D_t/2$, giving \eqref{eq:halfspace-map}.

The event $E_t=\{t<G<t+1/t,\ Z<-2\}$ lies in $A_t\setminus B_t$. Indeed $\sin x\ge5x/6$ for $0<x\le1$, and on this event
\[
 G\cos(1/t)+Z\sin(1/t)
 <t+1/t-5/(3t)<t.
\]
Monotonicity of $\varphi$ on $[t,\infty)$ gives
\[
 \Pr(t<G<t+1/t)\ge\frac{\varphi(t+1/t)}t
 \ge e^{-3/2}\frac{\varphi(t)}t.
\]
Furthermore
\[
 m_t\le\int_t^\infty(s/t)\varphi(s)\dd s=\varphi(t)/t.
\]
Independence of $G,Z$ therefore gives $\Pr(E_t)\ge c m_t$, proving the lower bound in \eqref{eq:halfspace-D}; its upper bound follows by the union bound. The same interval estimate is a lower bound for $m_t$ and proves \eqref{eq:mills}. The identity for $\delta_t$ follows from elementary trigonometry.
\end{proof}

\subsection{The logarithmic rare-event obstruction}
For $\beta>0$, choose $R_t>0$ uniquely by
\begin{equation}\label{eq:log-radius}
 m_tR_t^2[\log(e+R_t^2)]^\beta=1.
\end{equation}
The left side is continuous and strictly increasing in $R_t$ from zero to infinity. The targets in \Cref{lem:halfspaces} then satisfy \eqref{eq:log-class} with $M=1$ and equality.

We record the scale rigorously. Set $Q_t=R_t^2$ and $\ell_t=\log(1/m_t)$. From \eqref{eq:mills},
\begin{equation}\label{eq:tail-log}
 \ell_t=t^2/2+\log t+O(1)\sim t^2/2.
\end{equation}
Equation \eqref{eq:log-radius} first implies $Q_t\to\infty$, since a bounded subsequence would make its left side tend to zero. Taking logarithms gives
\[
 \ell_t=\log Q_t+\beta\log\log(e+Q_t).
\]
The second term is $o(\log Q_t)$, so $\log(e+Q_t)\sim\ell_t$. Returning to \eqref{eq:log-radius} yields
\begin{equation}\label{eq:log-radius-scale}
 m_tR_t^2=[\log(e+R_t^2)]^{-\beta}\sim(t^2/2)^{-\beta}.
\end{equation}
Write $w_t=W_2(\mu_t,\nu_t)$. The exact identities and bounds in \Cref{lem:halfspaces} now imply
\begin{equation}\label{eq:log-sharp-scales}
 w_t\sim2^{\beta/2}t^{-\beta-1},
 \qquad \norm{T_{\mu_t}-T_{\nu_t}}_2\asymp_\beta t^{-\beta}.
\end{equation}
For the second assertion, \eqref{eq:halfspace-map} is bounded below by $cm_tR_t^2$ and above by $(2+\delta_t^2)m_tR_t^2$. An exact asymptotic of $D_t/m_t$ is unnecessary.
Consequently, for every $\theta>\beta/(\beta+1)$,
\[
 \frac{\norm{T_{\mu_t}-T_{\nu_t}}_2}{w_t^\theta}
 \longrightarrow\infty.
\]
This gives the matching obstruction for $0<\beta\le1/2$.

At $\beta=0$, take $R_t=m_t^{-1/2}$. Both second moments are exactly one, while \eqref{eq:halfspace-W} gives $w_t=\delta_t\to0$ and \eqref{eq:halfspace-map}--\eqref{eq:halfspace-D} give a uniform positive lower bound for map distance. This excludes every uniform modulus at the endpoint, not just positive H\"older exponents.

\subsection{The bounded one-third obstruction}
For completeness we reconstruct the Gaussian three-atom family from \cite[Theorem~1.3]{ManuscriptV3}, which adapts the three-plane mechanism in \cite{OpenAI2026}.
Let
\[
 F(s)=\Pr(|G|<s)\quad(s\ge0),\qquad b=ka,
\]
where $k>0$ is fixed and $a\downarrow0$. There is a unique $c_a\in\R$ satisfying
\begin{equation}\label{eq:three-mass}
 \E F((c_a+kaZ)_+)=F(a)=:m_a.
\end{equation}
The expectation is continuous, strictly increasing in $c_a$, and ranges from zero to one. Define
\[
 u_a(x)=\max\{-x_1,x_1,a\},\qquad
 v_a(x)=\max\{-x_1,x_1,c_a+ka x_2\}.
\]
Their central cells are $A_a=\{|x_1|<a\}$ and
$B_a=\{|x_1|<(c_a+ka x_2)_+\}$. They have equal mass by \eqref{eq:three-mass}. Symmetry in $x_1$ gives the induced targets
\begin{align*}
 \mu_a&=\tfrac12(1-m_a)(\delta_{-e_1}+\delta_{e_1})+m_a\delta_0,\\
 \nu_a&=\tfrac12(1-m_a)(\delta_{-e_1}+\delta_{e_1})+m_a\delta_{be_2}.
\end{align*}
Their convex gradients are the optimal maps. Matching the two outer atoms and moving the central atom costs $m_ab^2$, and every coupling has at least that second-coordinate cost. Hence
\begin{equation}\label{eq:three-W}
 W_2(\mu_a,\nu_a)^2=m_ab^2.
\end{equation}

Let $c_*$ be the unique real number with $\E(c_*+kZ)_+=1$. Since $F(s)\le2\varphi(0)s$ and $F(as)/a\to2\varphi(0)s$, dominated convergence gives, for each fixed $c$,
\[
 a^{-1}\E F(a(c+kZ)_+)\longrightarrow2\varphi(0)\E(c+kZ)_+.
\]
Monotone bracketing at fixed $c_-<c_*<c_+$ in \eqref{eq:three-mass}, followed by arbitrarily close brackets, proves $c_a/a\to c_*$. In particular these ratios are bounded. The two central intervals are nested on every $x_2$ slice, so
\[
 D_a:=\gamma_d(A_a\mathbin\triangle B_a)
 =\E|F((c_a+kaZ)_+)-F(a)|.
\]
The integrand divided by $a$ is dominated by a fixed multiple of $1+|Z|$. Thus
\[
 D_a/a\longrightarrow2\varphi(0)\E|(c_*+kZ)_+-1|>0.
\]
The limit is positive because $(c_*+kZ)_+$ is not constant. Coordinate orthogonality gives exactly
\[
 |T_{\mu_a}-T_{\nu_a}|^2
 =\ind_{A_a\mathbin\triangle B_a}+b^2\ind_{B_a},
 \qquad \norm{T_{\mu_a}-T_{\nu_a}}_2^2=D_a+m_ab^2.
\]
Since $m_a\sim2\varphi(0)a$, the target distance is of order $a^{3/2}$ and the map distance of order $a^{1/2}$. All targets lie in $B_1$ for small $a$.

To fit any fixed logarithmic class with $M=1$, multiply every atom by
\[
 s_\beta=[\log(e+1)]^{-\beta/2}\le1.
\]
On the scaled support, $|y|^2[\log(e+|y|^2)]^\beta\le s_\beta^2[\log(e+1)]^\beta=1$. Scaling the convex potentials scales the maps, so both map distance and $W_2$ are multiplied by $s_\beta$. Every exponent above $1/3$ still fails. This is the matching obstruction for $\beta\ge1/2$; at the transition both constructions match.

Together with the upper proof, this completes \Cref{thm:log} in $d\ge2$. All examples use only two source coordinates. In $d=1$, the two increasing Brenier maps use a common source-quantile parameter, so their joint law is the monotone optimal coupling of the targets. This proves \eqref{eq:one-dimensional} and its stated exceptional scope.

\section{Uniform full-support smooth sources}\label{sec:smooth}
\begin{proof}[Proof of \Cref{thm:smooth}]
Strong convexity gives a unique minimizer $x_*$ of $V$ with $\nabla V(x_*)=0$. Put $S(x)=|x-x_*|$. Strong convexity and integration of the gradient-Lipschitz bound give
\[
 |\nabla V(x)|\le\Lambda S(x),\qquad
 V(x_*)+\tfrac\kappa2 S(x)^2\le V(x)
 \le V(x_*)+\tfrac\Lambda2 S(x)^2.
\]
Therefore $0<Z<\infty$ and
\begin{equation}\label{eq:smooth-envelope}
 r(x)\le(\Lambda/(2\pi))^{d/2}e^{-\kappa S(x)^2/2},
 \qquad \E_\rho e^{\kappa S^2/4}\le(2\Lambda/\kappa)^{d/2}.
\end{equation}
The normalization bound uses $Z\ge e^{-V(x_*)}(2\pi/\Lambda)^{d/2}$; the exponential-moment bound is the integral of the resulting Gaussian envelope.

Use the same $F(s)=\min(1,e^s)$ as above. The difference of minimum weights satisfies
\[
 \frac{m_{j,h}(x-he_j)-m_{j,h}(x)}{r(x)}
 =F(V(x)-V(x-he_j))-F(V(x)-V(x+he_j)).
\]
Because $F$ is $1$-Lipschitz with unit range, and $0<h\le1$,
\begin{align*}
 A_{j,h}(x)&\le\min\{|V(x+he_j)-V(x-he_j)|,1\}
 \le\min\{2h\Lambda(S(x)+1),1\},\\
 B_{j,h}(x)&\le\min\{|V(x+he_j)-V(x)|,1\}
 \le\min\{h\Lambda(S(x)+1),1\}.
\end{align*}
Indeed integration on the segment bounds the first potential difference by $2h\Lambda S+\Lambda h^2$, and the second by $h\Lambda S+\Lambda h^2/2$. Thus \Cref{lem:minimum} gives
\begin{equation}\label{eq:smooth-interpolation}
 \begin{split}
 \norm{\nabla u-\nabla v}_{L^2(\rho)}^2
 &\le12dh^{-2}\norm{u-v-a}_{L^2(\rho)}^2\\
 &\quad+K_\eta\Lambda^\eta h^\eta\sum_j\int
 (|\partial_j u|^2+|\partial_j v|^2)(1+S)^\eta\dd\rho.
 \end{split}
\end{equation}
Every signed-integral justification is already contained in the general minimum-weight lemma and uses only the original $L^2$ gradients.

For the centered Brenier potentials, apply \Cref{lem:pairing} with $A=|\partial_j u|^2/M^2$, the radial variable $S$, $a=\kappa/4$, and $H=(2\Lambda/\kappa)^{d/2}$. The pushforward identity and \eqref{eq:log-class} verify its target hypothesis, while \eqref{eq:smooth-envelope} verifies its source hypothesis. Its constant is precisely \eqref{eq:smooth-B}. Hence the sum of the weighted gradient integrals in \eqref{eq:smooth-interpolation} is at most $2dB_{d,\kappa,\Lambda,\beta}M^2$.

Since $W=V-\kappa|x|^2/2$ is finite convex, \Cref{thm:input} gives $\norm{u-v}_2^2\le(2/\kappa)w^2$. We obtain
\[
 \norm{T_\mu-T_\nu}_2^2
 \le(24d/\kappa)w^2h^{-2}
       +2dK_\eta\Lambda^\eta B_{d,\kappa,\Lambda,\beta}M^2h^\eta.
\]
For $0<w\le M$, take $h=(w/M)^{2/(\eta+2)}$. For $w>M$ use $\norm{T_\mu-T_\nu}_2\le2M$, and for $w=0$ use uniqueness. This proves \eqref{eq:smooth-main} with \eqref{eq:smooth-C}. Subtracting a common affine function gives the common-center statement.

For sharpness over the class, it contains the source $N(0,\kappa^{-1}I_d)$, whose gradient-Lipschitz constant is $\kappa\le\Lambda$. If $u$ is a potential for any standard-Gaussian example above, then $u(\sqrt\kappa x)/\sqrt\kappa$ has the same gradient law and map-discrepancy distribution under this dilated source. The target laws, logarithmic moments and Wasserstein distances are unchanged. The Gaussian obstructions therefore establish sharpness for the class in $d\ge2$. Individual sharpness for non-Gaussian members is not claimed.
\end{proof}

\section{The sharp critical order of finite-moment constants}\label{sec:q}
\begin{proof}[Proof of \Cref{thm:q}]
Set $\eta=1$ in \eqref{eq:gaussian-interpolation}, so $K_1=14$. H\"older's inequality with conjugate exponents $q/2$ and $p=q/(q-2)$ gives
\[
 \int|\partial_j u|^2(1+|x_j|)\dd\gamma_d
 \le D_q\norm{\partial_j u}_{L^q(\gamma_d)}^2.
\]
For centered Brenier potentials, \Cref{thm:input} and the gradient laws imply
\begin{equation}\label{eq:q-balance}
 \norm{T_\mu-T_\nu}_2^2
 \le24d w^2h^{-2}+14dD_qM_q^2h.
\end{equation}
Here the sum of coordinate $L^q$ norms squared is at most $d$ times the sum of the two vector $L^q$ norms squared, hence at most $dM_q^2$.
If $w>0$, then $M_q>0$, and $w\le M_q$ by the triangle inequality through $\delta_0$ and $L^q$ domination of $L^2$. Since $D_q\ge1$, the choice
\begin{equation}\label{eq:q-scale}
 h=\left(\frac{w}{M_q\sqrt{D_q}}\right)^{2/3}
\end{equation}
lies in $(0,1]$. Substitution in \eqref{eq:q-balance} gives coefficients $24d$ and $14d$ multiplying the common scale $D_q^{2/3}M_q^{4/3}w^{2/3}$. This proves \eqref{eq:q-upper}. The zero cases follow from uniqueness; $M_q=0$ means both targets are $\delta_0$.

For every $p\ge1$, maximizing $s^pe^{-s^2/4}$ gives
\[
 \E|G|^p\le(2p/e)^{p/2}\E e^{G^2/4}
 =\sqrt2(2p/e)^{p/2}.
\]
Thus
\[
 D_q\le1+2\sqrt{p/e}\le(1+2/\sqrt e)\sqrt p,
\]
and the upper bound in \eqref{eq:q-sharp-order} holds with
\begin{equation}\label{eq:q-Ad}
 A_d=\sqrt{38d}(1+2/\sqrt e)^{1/3}.
\end{equation}
Retaining $D_q$ in the admissible scale \eqref{eq:q-scale} is essential to this dependence; the scale need not be the exact numerical minimizer.

For the matching lower bound, fix $d\ge2$, $q>2$, and use \Cref{lem:halfspaces} with any $t\ge1$ and radius $R=m_t^{-1/q}$. Both $q$th moments are exactly one, so $M_q=2$. Since $1-2/q=1/p$, equations \eqref{eq:halfspace-W}--\eqref{eq:halfspace-D} imply
\begin{equation}\label{eq:q-ratio-lower}
 \frac{\norm{T_\mu-T_\nu}_2}{M_q^{2/3}W_2(\mu,\nu)^{1/3}}
 \ge2^{-2/3}\sqrt c\,m_t^{1/(3p)}\delta_t^{-1/3}.
\end{equation}
Take $t=\sqrt p\ge1$. The lower bound in \eqref{eq:mills} gives
\[
 m_t\ge e^{-3/2}(2\pi)^{-1/2}e^{-p/2}p^{-1/2}.
\]
Therefore, using $(\log p)/p\le1/e$,
\begin{align*}
 m_t^{1/(3p)}
 &\ge\exp\left(-\frac16-\frac1{2p}
                   -\frac{\log(2\pi)}{6p}-\frac{\log p}{6p}\right)\\
 &\ge\exp\left(-\frac16-\frac12-\frac{\log(2\pi)}6-\frac1{6e}\right).
\end{align*}
Also $\delta_t=2\sin(1/(2t))\le1/t=p^{-1/2}$. Inserting these bounds into \eqref{eq:q-ratio-lower} proves the lower inequality in \eqref{eq:q-sharp-order} with
\begin{equation}\label{eq:q-c0}
 c_0=2^{-2/3}\sqrt{e^{-3/2}\Pr(G<-2)}
 \exp\left(-\frac16-\frac12-\frac{\log(2\pi)}6-\frac1{6e}\right)>0.
\end{equation}
This holds for every $p>1$, not just asymptotically, and the source is fixed throughout. Since $p\sim2/(q-2)$ as $q\downarrow2$, the critical order follows.

In $d=1$, the quantile identity gives map distance $w$, and $w\le M_q$ implies $w\le M_q^{2/3}w^{1/3}$. Two opposite nonzero Dirac targets have $w=M_q$, so equality is attainable and $C_*(1,q)=1$.
\end{proof}

\section{Scope and relation to the earlier manuscript}
As an integrability requirement, the logarithmic condition is weaker than any fixed higher-power moment condition. Its positive exponent is nevertheless possible because the minimum-density identity produces a polynomial Gaussian-coordinate weight, and the logarithmic target moment pairs with Gaussian exponential-square integrability. Neither that pairing nor the proof assumes independence between transport values and source points.

Version~3 \cite{ManuscriptV3} contains additional results for compact log-concave sources and hard-boundary target-tail estimates. They remain separate, established parts of that earlier manuscript. The present focused version does not extend the logarithmic theorem to arbitrary strongly log-concave sources with hard boundaries, nor to arbitrary bounded non-log-concave perturbations. The smooth-source theorem requires full support and the stated global gradient-Lipschitz bound. Sharpness for that class follows through its Gaussian member; no claim is made for every individual non-Gaussian source.

The sharpness of the exponent in \Cref{thm:log}, the absence of a modulus at $\beta=0$, and the critical order in \Cref{thm:q} concern fixed source dimension $d\ge2$. Dimension one is governed by the exact quantile isometry. Mathematical proof and bibliographic priority are distinct issues; the latter is not asserted here.

\clearpage
\bibliographystyle{plain}
\bibliography{references}
\end{document}
